\subsection{The full Schwartz trace and arithmetic principal parts}
\label{sec:cq-principalparts}
Retain the source domain
\[
 \mathcal S_0^{\rm ev}
 =\{\psi\in\mathcal S(\mathbb R):
       \psi(-x)=\psi(x),\ \psi(0)=\widehat\psi(0)=0\},
 \qquad
 \widehat\psi(y)=\int_{\mathbb R}\psi(x)e^{-2\pi ixy}\,dx.
\]
Its full trace map takes values in entire functions:
\begin{equation}\label{eq:cq-Schwartz-trace}
 \begin{aligned}
 F_\psi(s)&=\mathcal F_\mu\mathcal E\psi(s),\\
 F_\psi(i(z-\tfrac12))&=\zeta(z)\mathcal M\psi(z),
 &\mathcal M\psi(z)&=\int_0^\infty\psi(x)x^z\,\frac{dx}{x}.
 \end{aligned}
\end{equation}
The Mellin integral is holomorphic for $\operatorname{Re}z>-2$,
and the factorization holds there with its removable value at
$z=1$. This statement does not impose an unproved order bound
on the entire $F_\psi$ for arbitrary Schwartz inputs.

Here is a direct proof with all analytic domains retained.
For fixed $x>0$, periodize $\psi$ by the smooth function
\[
 y\longmapsto\sum_{k\in\mathbb Z}\psi(x(k+y)).
\]
Its Fourier coefficients are $x^{-1}\widehat\psi(\ell/x)$:
integrate on $y\in[0,1]$, substitute $v=x(k+y)$, and sum the
adjacent intervals, justified by Schwartz decay. These coefficients
decrease faster than every power of $\ell$, so their Fourier series
converges with all derivatives and at $y=0$ gives the Poisson
identity. Evenness and the two exact zero conditions yield
\[
 \sum_{n\geq1}\psi(nx)=x^{-1}\sum_{n\geq1}\widehat\psi(n/x),
 \qquad \mathcal E\psi(x)=\mathcal E\widehat\psi(1/x).
\]
For $x\geq1$, a Schwartz bound of arbitrarily high order gives
$|\sum\psi(nx)|\leq C_Ax^{-A}\sum n^{-A}$ for $A>1$.
The same estimate after termwise differentiation follows by choosing
the Schwartz exponent larger than the derivative order plus one.
The Poisson identity supplies the corresponding bounds at zero.
Thus $\mathcal E\psi(e^u)$ decreases faster than $e^{-A|u|}$
for every $A>0$ at both ends. Its Fourier integral and each
derivative in $s$ are dominated locally uniformly in $s$ by an
integrable function. This proves that $F_\psi$ is entire.

Since $\psi$ is even and $\psi(0)=0$, its Taylor formula gives
$\psi(x)=O(x^2)$ at zero. Together with Schwartz decay at infinity
this proves the stated holomorphy of $\mathcal M\psi$ by
dominated differentiation, including all powers of $\log x$.
For $\sigma=\operatorname{Re}z>1$,
\[
 \sum_{n\geq1}\int_0^\infty|\psi(nx)|x^\sigma\,\frac{dx}{x}
 =\Bigl(\sum_{n\geq1}n^{-\sigma}\Bigr)
          \int_0^\infty|\psi(y)|y^\sigma\,\frac{dy}{y}<\infty.
\]
Absolute interchange and the original exponent $x^{-is}$ now give
\eqref{eq:cq-Schwartz-trace} on this half-plane. Moreover
$\mathcal M\psi(1)=\int_0^\infty\psi(x)\,dx
=\widehat\psi(0)/2=0$. This cancels the simple pole of $\zeta$
at $1$, so the right side is holomorphic for
$\operatorname{Re}z>-2$. The identity theorem extends the equality
to that entire half-plane, proving the assertion.

Let $\mathcal Z$ denote the set of distinct nontrivial zeros $\rho$
of $\zeta$, with their original multiplicities $m_\rho$.
The evenness of $\Xi$, immediate from the retained cosine integral,
gives the exact reflected-coordinate identity
\begin{equation}\label{eq:cq-reflected-Xi}
 s=i(z-\tfrac12),\qquad z=\tfrac12-is,\qquad
 \Xi(i(z-\tfrac12))=\xi(1-z)=\xi(z).
\end{equation}
At each $\rho\in\mathcal Z$ the completion factor
$z(z-1)\pi^{-z/2}\Gamma(z/2)/2$ is a nonvanishing holomorphic
unit. Hence $\lambda_\rho=i(\rho-\tfrac12)$ is a zero of
$\Xi$ of exactly multiplicity $m_\rho$, and these are all its
zeros. No reflection of the retained spectral operator $s$ has
been performed: \eqref{eq:cq-reflected-Xi} is the exact Mellin
parameter dictionary for that same operator. The factorization
\eqref{eq:cq-Schwartz-trace} proves
\begin{equation}\label{eq:cq-fulltrace-jets}
 \begin{gathered}
 D^jF_\psi(\lambda_\rho)=0
 \quad(\rho\in\mathcal Z,\ 0\leq j<m_\rho),\\
 F_\psi\in\cqE\ \Longrightarrow\ F_\psi\in\Xi\cqE.
 \end{gathered}
\end{equation}
The first assertion follows because $\mathcal M\psi$ is holomorphic
at every $\rho$ and $s-\lambda_\rho=i(z-\rho)$; the second
then follows from the proved division theorem. It is an exact
intersection assertion about the full Schwartz trace, without
enlarging the order-at-most-one domain or closing it in $\tau$.

For each $\rho$ retain the actual local unit
$a_\rho(z)=\zeta(z)/(z-\rho)^{m_\rho}$ and the space
\[
 \mathcal P_\rho^-=
  \left\{\sum_{j=1}^{m_\rho}c_{-j}(z-\rho)^{-j}:c_{-j}\in\cqC\right\}.
\]
The principal-part map on the original entire-function set is
\begin{equation}\label{eq:cq-principalparts}
 \mathfrak P:\cqE\longrightarrow\prod_{\rho\in\mathcal Z}\mathcal P_\rho^-,
 \qquad
 F\longmapsto\left(\operatorname{PP}_{z=\rho}
                 \frac{F(i(z-\tfrac12))}{\zeta(z)}\right)_\rho.
\end{equation}
The kernel of this map is exactly $\Xi\cqE$. It induces an
injective linear map from the underlying vector space of each
$Q_\gamma,Q_\kappa$ onto the same specified image. Every finite
projection is surjective. The map from $Q_\kappa$ to its image,
with the product subspace topology, is a topological linear
isomorphism; the product of all principal-part spaces is its
completion. From $Q_\gamma$ it is a continuous bijection onto
that image whose inverse is not continuous.

To prove each claim while retaining the analytic units, write
$a_\rho(\rho+w)^{-1}=\sum_{\ell\geq0}\alpha_{\rho,\ell}w^\ell$.
The exact Laurent coefficient at degree $-m_\rho+q$, $0\leq q<m_\rho$,
of the $\rho$ component is
\begin{equation}\label{eq:cq-principalpart-coefficients}
 c_{-m_\rho+q}
 =\sum_{j=0}^{q}\frac{i^j}{j!}F^{(j)}(\lambda_\rho)
                                  \alpha_{\rho,q-j}.
\end{equation}
This follows by multiplying the two convergent local Taylor
series and the retained factor $w^{-m_\rho}$. Its finite matrix
is triangular with nonzero diagonal entries
$i^j/(j!a_\rho(\rho))$. Explicitly, given a prescribed principal
part $P_\rho(w)$, take the coefficients $b_j$ through degree
$m_\rho-1$ of $a_\rho(\rho+w)w^{m_\rho}P_\rho(w)$;
the required original derivatives are $j!i^{-j}b_j$.
Multiplying back by $a_\rho^{-1}$ recovers the prescribed negative
powers, since the discarded error is divisible by $w^{m_\rho}$.
This proves a finite-dimensional linear isomorphism at every zero,
including its inverse and every factorial and power of $i$.

The vanishing of all principal parts is therefore equivalent to
the vanishing of all the original multiplicity jets of $F$.
Theorem~\ref{thm:cq-equalclosures} identifies that kernel exactly
with $\Xi\cqE$. At any finite set of zeros, the required derivatives
can be realized by the explicit Hermite lift
\eqref{eq:cq-finite-Hermite-lift}; this proves surjectivity of
every finite principal-part projection. Each block map and its
inverse are continuous because their finite matrices have fixed
complex coefficients. Their product is therefore a homeomorphism
of the two product spaces: the inverse image of a cylinder fixing
finitely many coordinates is again controlled by those finitely
many blocks. Applying this product isomorphism to
Theorem~\ref{thm:cq-jet-topology} proves all stated topology,
completion, and comparison assertions for $\mathfrak P$.

This map also carries the original spectral action exactly. For
$P_\rho=\mathfrak P_\rho(F)$ one has
\begin{equation}\label{eq:cq-s-principalparts}
 \mathfrak P_\rho(sF)
 =\operatorname{PP}_{w=0}
                  ((\lambda_\rho+iw)P_\rho(w)).
\end{equation}
Indeed the full Laurent quotient for $sF$ is the quotient for $F$
multiplied by $i(z-\tfrac12)=\lambda_\rho+iw$.
Its holomorphic residual contributes no negative power after this
multiplication. Thus the finite block is the retained
$\lambda_\rho$ plus the explicitly specified nilpotent operator
$P\mapsto\operatorname{PP}(iwP)$, whose nilpotency index is
$m_\rho$: iteration on $w^{-m_\rho}$ gives a nonzero multiple
of $w^{-1}$ at step $m_\rho-1$ and zero at step $m_\rho$.

For $I=\Xi\cqE$, the derivative ambiguity
$W_k=(D^kI+I)/I$ has, under $\mathfrak P$, the exact finite image
\begin{equation}\label{eq:cq-principal-ambiguity}
 \prod_{\rho\in F}
 \left\{\sum_{j=1}^{\min(k,m_\rho)}c_{-j}(z-\rho)^{-j}:c_{-j}\in\cqC\right\}
 \quad(F\subset\mathcal Z\text{ finite}).
\end{equation}
The finite image in \eqref{eq:cq-ambiguityjets} consists exactly
of numerator jets divisible by $w^{\max(m_\rho-k,0)}$ after
the invertible substitution $s-\lambda_\rho=iw$.
Multiplication by the nonvanishing unit $a_\rho^{-1}$ preserves
this divisibility and is invertible on the truncated Taylor ring.
Division by $w^{m_\rho}$ therefore yields exactly all negative
powers of order at most $\min(k,m_\rho)$, proving both
containment and surjectivity in the displayed finite image.
In particular the complete local-unit expression, rather than only
a count of vanishing orders, transports the derivative defect to
the arithmetic Mellin poles.
